\section{Introduction}

True random number generators (TRNGs) are essential building blocks in modern secure systems, providing the unpredictable seeds and nonces required by cryptographic protocols. Unlike pseudorandom number generators, TRNGs derive randomness from physical entropy sources whose unpredictability can be grounded in device physics rather than computational assumptions~\cite{stipcevic2014true}. However, dedicated entropy source circuits (ex. ring oscillator jitter harvesters~\cite{sunar2007provably} or metastable latch designs~\cite{tokunaga2008true}) occupy silicon area and consume power with no secondary function, making them unattractive for area- and energy-constrained embedded platforms.

This inefficiency has motivated the study of \emph{entropy harvesting}, which repurposes existing on-chip structures as dual-use entropy sources. SRAM power-up states~\cite{holcomb2009power}, DRAM retention noise~\cite{kim2019d}, and Flash threshold voltage variation~\cite{wang2012flash} have each been exploited in this spirit, amortizing the cost of entropy generation across components that already serve a storage or compute function.

Resistive RAM (ReRAM) is a compelling candidate for embedded entropy harvesting. Its stochastic filament formation and rupture dynamics produce measurable cycle-to-cycle variation in switching behavior~\cite{yu2016resistive, balatti2015true}, and its increasing adoption in low-power microcontrollers and edge devices~\cite{ramxeedreramproducts, nuvoton_rram} make ReRAM switching entropy increasingly available alongside it's nonvolatile storage usage. Several prior studies have proposed ReRAM-based TRNGs that exploit write latency or resistance variation as entropy observables~\cite{huang2012contact, jiang2017novel, wei2016true, pang2019memristors, tolga2024trng}.

Yet these prior approaches share a critical blind spot. Recent work demonstrates that the existing standalone commercial ReRAM memories currently on the market implement a linear block error-correcting code (ECC) \cite{sheaves2026brewrc}. ECCs introduce additional parity bits to the datawords programmed by users, forming codewords. Without knowledge of the ECC, applied datapatterns inject unexpected, deterministic, switching activity into every write transaction. This observation allows the authors of \cite{sheaves2026brewrc} to reverse engineer the bit-exact ECC. Prior evaluations on commercial-off-the-shelf ReRAM entropy sources applied random data patterns to random addresses and aggressive debiasing techniques making the true entropy of these devices unclear. This practice both injects structural bias from the ECC and violates the stationarity assumptions of standard entropy assessments (i.e. NIST SP~800-90B)~\cite{sonmez2016recommendation}.

In this paper, we present a novel technique to use knowledge of reverse engineered ECC embedded in two commercial COTS ReRAM devices to characterize and leverage the stochastic, \emph{codeword-level} switching behavior of ReRAM crossbar memories, allowing for the construction of a far more efficient entropy source. We make three key contributions: (1)~we show that ECC-aware analysis explains structural biases that have weakened prior ReRAM entropy claims; (2)~we isolate and vet dedicated memory regions, applying a wear-leveled burst sequence whose Hamming distance profile is optimized to maximize the min-entropy of the corresponding write latency; (3)~we estimate $H_{\min}$ of the vetted write patterns/regions using NIST SP~800-90B and validate a constructed TRNG against both the NIST SP~800-22 statistical test suite~\cite{rukhin2001statistical} and PractRand~\cite{practrand} on two commercially available ReRAM modules. 

Our results demonstrate a TRNG achieving a throughput of 1.7~kbit/s on a COTS MB85AS8MT 8Mb ReRAM device from Ramxeed.